% SPDX-License-Identifier: LPPL-1.3c
% Copyright (C) 2026 Christophe Jorssen
% Author and Current Maintainer: Christophe Jorssen <christophe.jorssen@gmail.com>
% LPPL maintenance status: maintained
% This file is part of luafemm. See LICENSE and LICENSES.md for its terms.
\documentclass[varwidth=165mm,border=4mm]{standalone}
\usepackage{tikz,pgfplots}
\usetikzlibrary{femm}
\usepgfplotslibrary{femm}
\pgfplotsset{compat=1.18}
\definecolor{ndtblue}{RGB}{25,83,112}
\definecolor{ndtred}{RGB}{192,69,44}
\def\Mesh{.4}\def\Frequency{2000}
\newbox\NDTCalculation
\def\ShallowR{}\def\ShallowX{}\def\ShallowPlane{}
\def\DeepR{}\def\DeepX{}\def\DeepPlane{}
\tikzset{ndt/model/.style={femm/problem={type=axisymmetric,
    frequency=\Frequency,xmax=45,ymin=-45,ymax=45,mesh size=2},
    femm/circuit={name=probe,current=1},
    femm/materials/specimen/.style={mu r=1,conductivity=58000000}}}
\pgfplotsset{ndt/axis/.style={grid=major,grid style={gray!15},
    axis line style={gray!60},tick label style={font=\small},
    legend style={draw=none,fill=white,font=\small},
    scaled ticks=false},
  ndt/shallow/.style={ndtblue,thick,mark=*,mark size=1.3pt},
  ndt/deep/.style={ndtred,thick,mark=square*,mark size=1.5pt}}
\pgfmathsetmacro{\NearMesh}{1.5*\Mesh}
\pgfmathsetmacro{\GrooveMesh}{\Mesh/2}
\begin{document}
\centering
% Each position is an independent steady harmonic calculation. The coil
% moves between solves: there is no velocity or motional-current term.
\foreach \Position [count=\Index] in {-12,-10,-8,-6,-5,-4,-3,-2,-1,0,1,2,3,4,5,6,8,10,12}{%
  \foreach \Case/\Depth in {sound/0,shallow/.5,deep/1.5}{%
    % Discard the picture box, retaining the named numerical measurement.
    \setbox\NDTCalculation=\hbox{%
      \tikzpicture[ndt/model,
          femm/cache={file=\jobname-\Case-\Index,mode=auto}]
        \path[femm/region={material=air,mesh size=\NearMesh}]
          (0,-15) rectangle (12,15);
        \path[femm/region={material=specimen,mesh size=\Mesh}]
          (5,-25) rectangle (8,25);
        \ifdim\Depth pt>0pt
          \path[femm/region={material=air,mesh size=\GrooveMesh}]
            ({8-\Depth},-1.5) rectangle (8,1.5);
        \fi
        \path[femm/region={material=copper,circuit=probe,turns=60,
            winding fill=.6,mesh size=\Mesh}]
          (9,{\Position-1.5}) rectangle (11,{\Position+1.5});
        \femmmeasure[name=\Case-\Index,quantity=circuit,circuit=probe,
          parameter=\Position]
      \endtikzpicture}%
  }
  % Use PGF's floating-point unit for differences of nearby impedances.
  % Convert ohms to milliohms only after subtracting the matched reference.
  \begingroup\pgfkeys{/pgf/fpu=true,/pgf/fpu/output format=fixed}
    \foreach \Case/\Prefix in {shallow/Shallow,deep/Deep}{%
      \pgfmathparse{1000*(\femmresult{\Case-\Index}{impedance_real}
        -\femmresult{sound-\Index}{impedance_real})}
      \edef\RealSignal{\pgfmathresult}
      \expandafter\xdef\csname\Prefix R\endcsname{%
        \csname\Prefix R\endcsname (\Position,\RealSignal)}
      \pgfmathparse{1000*(\femmresult{\Case-\Index}{impedance_imaginary}
        -\femmresult{sound-\Index}{impedance_imaginary})}
      \expandafter\xdef\csname\Prefix X\endcsname{%
        \csname\Prefix X\endcsname (\Position,\pgfmathresult)}
      \expandafter\xdef\csname\Prefix Plane\endcsname{%
        \csname\Prefix Plane\endcsname (\RealSignal,\pgfmathresult)}
    }
  \endgroup
}
\begin{tikzpicture}
  \begin{axis}[ndt/axis,width=8cm,height=6cm,xmin=-12,xmax=12,
      xlabel={Probe centre $s$ (mm)},ylabel={$\Delta R$ (m$\Omega$)},
      title={Resistive channel},legend pos=south east]
    \addplot[gray,densely dashed,forget plot] coordinates {(-12,0)(12,0)};
    \addplot[ndt/shallow] coordinates {\ShallowR};
    \addlegendentry{$d=0.5$ mm}
    \addplot[ndt/deep] coordinates {\DeepR};
    \addlegendentry{$d=1.5$ mm}
  \end{axis}
\end{tikzpicture}%
\begin{tikzpicture}
  \begin{axis}[ndt/axis,width=8cm,height=6cm,xmin=-12,xmax=12,
      xlabel={Probe centre $s$ (mm)},ylabel={$\Delta X$ (m$\Omega$)},
      title={Reactive channel}]
    \addplot[gray,densely dashed,forget plot] coordinates {(-12,0)(12,0)};
    \addplot[ndt/shallow] coordinates {\ShallowX};
    \addplot[ndt/deep] coordinates {\DeepX};
  \end{axis}
\end{tikzpicture}
\par\medskip
\begin{tikzpicture}
  \begin{axis}[ndt/axis,width=10cm,height=7cm,
      xlabel={$\Delta R$ (m$\Omega$)},ylabel={$\Delta X$ (m$\Omega$)},
      title={The same scan in the impedance plane},legend pos=north east]
    \addplot[ndt/shallow] coordinates {\ShallowPlane};
    \addlegendentry{$d=0.5$ mm}
    \addplot[ndt/deep] coordinates {\DeepPlane};
    \addlegendentry{$d=1.5$ mm}
    \node[anchor=north east,font=\small] at (axis cs:0,0) {Far from groove};
  \end{axis}
\end{tikzpicture}
\par\smallskip\small
$\Delta Z(s)=Z_{\rm groove}(s)-Z_{\rm sound}(s)$;
$f=2$ kHz, 60 turns, $\widehat I=1$ A, groove width 3 mm.
Markers are solved positions; straight segments only guide the eye.
The symmetric passage nearly retraces its impedance-plane trajectory.
\end{document}
